\documentclass[11pt]{article}

\usepackage[margin=0.85in]{geometry}
\usepackage[hidelinks]{hyperref}
\usepackage{enumitem}
\usepackage{titlesec}

\pagestyle{empty}
\setlength{\parskip}{0.55em}
\setlength{\parindent}{0pt}

\titleformat{\paragraph}[runin]{\bfseries}{}{}{}[.]
\titlespacing*{\paragraph}{0pt}{0.8em}{0.6em}

\begin{document}

\begin{center}
{\Large\textbf{Fitting Person Detection on a Plain Microcontroller}}\\[0.3em]
Akash Reddy K \quad\textbar\quad Independent research / portfolio project\\
Code, data, full paper: \url{https://github.com/akash-reddy-k/edgeAi}
\end{center}

\vspace{0.4em}

Commercial CCTV analytics (Hikvision AcuSense, Ambarella CV7) run person
detection on purpose-built AI silicon with dedicated neural accelerators.
This project asks how far equivalent functionality---loitering detection on
a fixed-mount camera---can be pushed onto a \textbf{plain microcontroller
with no such accelerator}, and measures the answer rather than assuming it.

\paragraph{Method} Fine-tune YOLOv8n for single-class person detection
(ShanghaiTech Campus dataset), track with ByteTrack, export to ONNX, apply
INT8 post-training quantization. An ONNX-graph profiler then computes flash
and activation-RAM footprint directly---no vendor toolchain or account
required, fully reproducible from committed code.

\paragraph{Finding 1 --- quantization alone is not enough} The INT8 model
is 72\% smaller than its FP32 export for a 0.2\% accuracy cost, but at the
stock YOLOv8n architecture, \textbf{nothing fits any plain STM32, at any
input resolution}---flash alone overshoots the largest plain part
(STM32H743, 2\,MB) by 50\%.

\paragraph{Finding 2 --- that is an architecture fact, not a hardware fact}
Profiling shows YOLOv8n's large-object detection pathway (P5/32) accounts
for 51.8\% of parameters and is functionally unused in fixed-mount overhead
CCTV, where a person never fills half the frame. Removing it cuts flash
2.8$\times$---more than quantization achieved on its own.

\paragraph{Finding 3 --- the fix does not cost accuracy, it improves it}
Reinvesting that freed budget as extra width (rather than shrinking every
scale uniformly) wins on every metric in a controlled, random-init-matched
comparison, including $+3.3$ points of mAP50-95. The resulting architecture
fits an STM32H743 at $320\times320$ with roughly 50\% headroom on both
flash and RAM---no NPU, no external memory.

\paragraph{Honest caveat} The accuracy numbers above are a relative
comparison on a held-constant, self-referential label set (no
ground-truth bounding-box annotations exist for this dataset); the MCU fit
numbers are an analytical lower bound, not yet verified on physical
hardware. Both limitations are stated explicitly in the full paper, and
both are active next steps.

\paragraph{Why this, not a bigger model} The research contribution is the
compression-and-deployment pipeline and its diagnostic method---distinguish
what a flash budget is actually spent on before assuming quantization is
the only lever---not a new detector.

\vspace{0.3em}
\noindent Full paper (12 pages, arXiv-style): \texttt{writeup/paper.pdf} in the repo above.

\end{document}
